% Drawing helpers of this talk, loaded after the theme.
\usetikzlibrary{decorations.pathreplacing}% braces

% Plot data: the 25 measured alloys, their front, the processing design, values read off the paper's figures,
% and the Euler colours of the colour slide; written by tools/make_assets.py.
% Written by tools/make_assets.py; edit the script, not this file.
% The 25 measured alloys (assets/liao2026/dataset.csv): yield strength (MPa)/elongation (%)/ID.
\def\alloydata{470.23/69.32/1,468.47/68.7/2,434.94/72.51/3,394.53/78.29/4,362.26/82.18/5,377.75/78.15/6,399.05/75.03/7,613.19/53.39/8,700.46/49.44/9,827.46/40.76/10,897.42/35.42/11,793.37/40.48/12,594.08/50.53/13,441.99/69.56/14,395.19/74.56/15,429.82/67.24/16,451.36/66.3/17,593.98/45.55/18,790.98/39.15/19,983.38/33.63/20,938.12/32.69/21,844.78/36.77/22,807.42/35.94/23,462.32/64.45/24,373.67/70.09/25}
% Their front: \drawfront[options] draws its staircase in (MPa,%) coordinates.
\newcommand\drawfront[1][]{\draw[#1] (362.26,82.18) -- (362.26,78.29) -- (394.53,78.29) -- (394.53,75.03) -- (399.05,75.03) -- (399.05,72.51) -- (434.94,72.51) -- (434.94,69.56) -- (441.99,69.56) -- (441.99,69.32) -- (470.23,69.32) -- (470.23,53.39) -- (613.19,53.39) -- (613.19,49.44) -- (700.46,49.44) -- (700.46,40.76) -- (827.46,40.76) -- (827.46,36.77) -- (844.78,36.77) -- (844.78,35.42) -- (897.42,35.42) -- (897.42,33.63) -- (983.38,33.63);}
% The processing design: hot-rolling reduction (%)/annealing temperature (C)/annealing time (h)/yield strength (MPa)/shade of blue (%).
\def\alloygrid{0/650/1/470/19,0/750/5/468/19,0/850/10/435/15,0/950/15/395/11,0/1050/20/362/8,20/1050/15/378/10,20/950/10/399/12,20/850/5/613/33,20/750/1/700/42,20/650/20/827/54,30/650/15/897/61,30/750/10/793/51,30/850/20/594/31,30/950/1/442/16,30/1050/5/395/11,40/1050/1/430/15,40/950/5/451/17,40/850/15/594/31,40/750/20/791/51,40/650/10/983/70,50/650/5/938/65,50/750/15/845/56,50/850/1/807/52,50/950/20/462/18,50/1050/10/374/9}
% Read off the paper's figures: the five candidates of Fig. 4f, yield strength/elongation/label (0: none),
% and candidate 1 in Fig. 5h as simulated and as measured on the processed alloy.
\def\candidatedata{556/68.8/0,666/67.6/0,677/65.3/2,823/48.1/1,899/46.5/0}
\def\predictedpoint{823,48.1}
\def\simulatedpoint{808,44.5}
\def\measuredpoint{825,43.5}
% Euler colours of nearly equal orientations: phi1 = 2 and 358 degrees (Phi 30, phi2 45), and
% phi2 = 1 and 89 degrees (phi1 120, Phi 30), which a 90-degree cubic symmetry relates.
\definecolor{wrapA}{RGB}{1,85,128}
\definecolor{wrapB}{RGB}{254,85,128}
\definecolor{symA}{RGB}{85,85,3}
\definecolor{symB}{RGB}{85,85,252}
% The 25 processing routes, rows 1-13 and 14-25: ID & C & h & % & MPa & %.
\def\alloytableA{1 & 650 & 1 & 0 & 470 & 69.3\\2 & 750 & 5 & 0 & 468 & 68.7\\3 & 850 & 10 & 0 & 435 & 72.5\\4 & 950 & 15 & 0 & 395 & 78.3\\5 & 1050 & 20 & 0 & 362 & 82.2\\6 & 1050 & 15 & 20 & 378 & 78.2\\7 & 950 & 10 & 20 & 399 & 75.0\\8 & 850 & 5 & 20 & 613 & 53.4\\9 & 750 & 1 & 20 & 700 & 49.4\\10 & 650 & 20 & 20 & 827 & 40.8\\11 & 650 & 15 & 30 & 897 & 35.4\\12 & 750 & 10 & 30 & 793 & 40.5\\13 & 850 & 20 & 30 & 594 & 50.5\\}
\def\alloytableB{14 & 950 & 1 & 30 & 442 & 69.6\\15 & 1050 & 5 & 30 & 395 & 74.6\\16 & 1050 & 1 & 40 & 430 & 67.2\\17 & 950 & 5 & 40 & 451 & 66.3\\18 & 850 & 15 & 40 & 594 & 45.5\\19 & 750 & 20 & 40 & 791 & 39.1\\20 & 650 & 10 & 40 & 983 & 33.6\\21 & 650 & 5 & 50 & 938 & 32.7\\22 & 750 & 15 & 50 & 845 & 36.8\\23 & 850 & 1 & 50 & 807 & 35.9\\24 & 950 & 20 & 50 & 462 & 64.5\\25 & 1050 & 10 & 50 & 374 & 70.1\\}
% Non-dominated sorting of a toy trade-off (pareto_toy): points x/y/rank in the unit square, both
% objectives maximised; staircases of fronts 1 and 2; the crowding box of the middle point of front 1.
\def\paretopoints{0.133/0.950/1,0.334/0.396/4,0.186/0.692/3,0.633/0.420/2,0.891/0.189/1,0.113/0.884/2,0.634/0.566/1,0.207/0.612/3,0.754/0.333/2,0.628/0.398/3,0.856/0.055/2,0.800/0.146/2,0.778/0.393/1,0.525/0.900/1,0.555/0.482/2,0.386/0.788/2}
\newcommand\drawfrontone[1][]{\draw[#1] (0.133,0.950) -- (0.133,0.900) -- (0.525,0.900) -- (0.525,0.566) -- (0.634,0.566) -- (0.634,0.393) -- (0.778,0.393) -- (0.778,0.189) -- (0.891,0.189);}
\newcommand\drawfronttwo[1][]{\draw[#1] (0.113,0.884) -- (0.113,0.788) -- (0.386,0.788) -- (0.386,0.482) -- (0.555,0.482) -- (0.555,0.420) -- (0.633,0.420) -- (0.633,0.333) -- (0.754,0.333) -- (0.754,0.146) -- (0.800,0.146) -- (0.800,0.055) -- (0.856,0.055);}
\newcommand\drawcrowdbox[1][]{\draw[#1] (0.525,0.393) rectangle (0.778,0.900);}
% SSIM of the illustrative patch against changed copies (ssim_examples): \ssim<copy><L, C, S, Q for
% their product, M for the mean absolute difference>, over the whole patch on grey values
\newcommand\ssimBrightL{0.89}
\newcommand\ssimBrightC{1.00}
\newcommand\ssimBrightS{1.00}
\newcommand\ssimBrightQ{0.89}
\newcommand\ssimBrightM{0.24}
\newcommand\ssimFlatL{1.00}
\newcommand\ssimFlatC{0.70}
\newcommand\ssimFlatS{1.00}
\newcommand\ssimFlatQ{0.70}
\newcommand\ssimFlatM{0.08}
\newcommand\ssimShuffledL{1.00}
\newcommand\ssimShuffledC{1.00}
\newcommand\ssimShuffledS{0.02}
\newcommand\ssimShuffledQ{0.02}
\newcommand\ssimShuffledM{0.20}
\newcommand\ssimBlurredL{1.00}
\newcommand\ssimBlurredC{0.97}
\newcommand\ssimBlurredS{0.86}
\newcommand\ssimBlurredQ{0.84}
\newcommand\ssimBlurredM{0.05}


% ---------------------------------------------------------------- the paper's figures
% The figures of Liao et al. (2026) as published (CC BY-NC-ND 4.0). A slide shows a part of one figure,
% unmodified, and draws its own labels and marks beside it, never over it.
% \figpart{x}{y}{figure}{x0}{y0}{x1}{y1}{width}: in a tikzpicture[x=1mm,y=1mm], the pixel box (x0,y0)--(x1,y1)
% of figure-0<figure>.jpg (2000 pixels wide, origin at its top-left corner), scaled to <width> mm, with its
% top-left corner at (x,y). Afterwards (fp nw) and (fp se) are its corners, and \fpx{px} and \fpy{py} give
% the slide position (mm) of an image pixel, to place a label beside a panel.
\newcommand\figpart[8]{%
  \pgfmathsetmacro\fpscale{#8/(#6-#4)}%
  \pgfmathsetmacro\fpfull{2000*\fpscale}%
  \pgfmathsetmacro\fpheight{(#7-#5)*\fpscale}%
  \def\fpX{#1}\def\fpY{#2}\def\fpXo{#4}\def\fpYo{#5}%
  \coordinate (fp nw) at (#1,#2);
  \coordinate (fp se) at ({#1+#8},{#2-\fpheight});
  \begin{scope}
    \clip (fp nw) rectangle (fp se);
    \node[anchor=north west,inner sep=0pt] at ({#1-#4*\fpscale},{#2+#5*\fpscale})
      {\includegraphics[width=\fpfull mm]{assets/liao2026/figure-0#3.jpg}};
  \end{scope}}
\newcommand\fpx[1]{(\fpX+(#1-\fpXo)*\fpscale)}
\newcommand\fpy[1]{(\fpY-(#1-\fpYo)*\fpscale)}

% ------------------------------------------------------------------ the property plane
% Yield strength against elongation, drawn alike wherever it returns: the measured alloys as grey dots,
% their front as a grey dashed staircase, generated candidates as red squares.
% \planeaxes{xmin}{xmax}{ymin}{ymax}{width}{height}{x ticks}{y ticks} draws the axes of a plot <width> by
% <height> mm whose corner (xmin,ymin) is the current origin; a scope with the style plane then takes
% (MPa,%) coordinates.
\tikzset{
  alloy/.style={circle,fill=udsgreymid,inner sep=0pt,minimum size=1.7mm},
  candidate/.style={rectangle,fill=udsred,inner sep=0pt,minimum size=1.9mm},
  tick label/.style={font=\scriptsize,text=udsgrey,inner sep=0pt},
  axis label/.style={font=\footnotesize,text=udsdark,inner sep=0pt},
  front/.style={udsgreymid,line width=0.6pt,dash pattern=on 1.6pt off 1.2pt},
}
\newcommand\planeaxes[8]{%
  \pgfmathsetmacro\plsx{#5/(#2-#1)}\pgfmathsetmacro\plsy{#6/(#4-#3)}%
  \tikzset{plane/.estyle={x=\plsx mm,y=\plsy mm,shift={(-#1,-#3)}}}%
  \draw[udsgreymid,line width=0.5pt] (0,#6) -- (0,0) -- (#5,0);
  \begin{scope}[plane]
    \foreach \t in {#7} {
      \draw[udsgreymid,line width=0.5pt] (\t,#3) -- ++(0pt,-1mm);
      \node[tick label,anchor=north] at ([yshift=-1.7mm]\t,#3) {\t};}
    \foreach \t in {#8} {
      \draw[udsgreymid,line width=0.5pt] (#1,\t) -- ++(-1mm,0pt);
      \node[tick label,anchor=east] at ([xshift=-1.7mm]#1,\t) {\t};}
  \end{scope}}
% inside a scope with the style plane
\newcommand\drawalloys[1][]{\foreach \s/\e/\i in \alloydata {\node[alloy,#1] at (\s,\e) {};}}
% candidate k is the node (cand k) for the labelled candidates 1 and 2
\newcommand\drawcandidates[1][]{%
  \foreach \s/\e/\l in \candidatedata {\node[candidate,#1] (cand\l) at (\s,\e) {};}}
% a candidate's label, ① or ②, beside its square
% (its digit on the baseline of the text around it; anchor=base also overrides the anchor that a
% nested picture inherits from the node it is set in)
\newcommand\candlabel[1]{\tikz[baseline=(cl.base)]\node[anchor=base,circle,draw=udsred,line width=0.5pt,
  text=udsred,font=\fontsize{6.5}{7}\selectfont\bfseries,inner sep=0pt,minimum size=3.1mm] (cl) {#1};}

% ------------------------------------------------------------------------ strips
% A strip across the text width with its top at y (mm): red with a warning sign for a catch, green with a
% check mark for what holds.
\newcommand\warnstrip[2]{%
  \fill[redtint] (0,#1) rectangle (144,#1-5.6);
  \fill[udsred] (0,#1) rectangle (1.2,#1-5.6);
  \node[anchor=west,inner sep=0pt,font=\footnotesize,text=udsdark] at (3.6,#1-2.8)
    {\textcolor{udsred}{\faExclamationTriangle}\enspace #2};}
\newcommand\goodstrip[2]{%
  \fill[greentint] (0,#1) rectangle (144,#1-5.6);
  \fill[testgreen] (0,#1) rectangle (1.2,#1-5.6);
  \node[anchor=west,inner sep=0pt,font=\footnotesize,text=udsdark] at (3.6,#1-2.8)
    {\textcolor{testgreen}{\faCheck}\enspace #2};}

% ------------------------------------------------------------------------- cards
% \cardhead{x}{y}{width}{colour}{text}: a card's head, 5.4 mm tall, with its label on one baseline
\tikzset{card head/.style={anchor=north west,inner xsep=1.6mm,inner ysep=0pt,minimum height=5.4mm,
  text height=1.9mm,text depth=0.6mm,font=\footnotesize\bfseries,text=white,align=left}}
\newcommand\cardhead[5]{\pgfmathsetmacro\chw{#3-3.2}%
  \node[card head,fill=#4,minimum width=#3mm,text width=\chw mm] at (#1,#2) {#5};}
% lines of a card body, 7 or 8 pt, ragged right without hyphenation
\tikzset{card text/.style={anchor=north west,inner sep=0pt,font=\footnotesize,text=udsdark,align=left,
  execute at begin node={\hyphenpenalty=10000\raggedright}}}
% \cardbullet{text}: one line of a card with the theme's square bullet and a hanging indent
\newcommand\cardbullet[1]{\par\hangindent=3.4mm\hangafter=1\noindent
  \makebox[3.4mm][l]{\textcolor{udsblue}{\rule[0.3ex]{1.05mm}{1.05mm}}}#1\par\vspace{0.8mm}}

% ------------------------------------------------------------------------ labels
\tikzset{
  label text/.style={inner sep=0pt,font=\footnotesize,text=udsdark,align=left,
    execute at begin node={\hyphenpenalty=10000\raggedright}},
  small label/.style={inner sep=0pt,font=\scriptsize,text=udsgrey,align=left,
    execute at begin node={\hyphenpenalty=10000\raggedright}},
  icon frame/.style={draw=udsdark,line width=0.4pt},
}

% ------------------------------------------------------------------ glyphs
% \fixarrow{x0}{x1}{y}{label}: a broad dark arrow from (x0,y) to (x1,y), its label above it
\newcommand\fixarrow[4]{%
  \fill[udsdark] (#1,#3+1.5) -- (#2-3.4,#3+1.5) -- (#2-3.4,#3+3.4) -- (#2,#3) -- (#2-3.4,#3-3.4)
    -- (#2-3.4,#3-1.5) -- (#1,#3-1.5) -- cycle;
  \node[anchor=south,small label,text=udsdark,align=center,text width=24mm] at ({(#1+#2)/2},#3+4.4) {#4};}
% \descstrip{x}{y}{w}{h}: a descriptor of six numbers as six cells, the three of the mean blue and the three
% of the spread dark; (x,y) is the top-left corner, w the width of a cell, h its height
\newcommand\descstrip[4]{%
  \foreach \dsj in {0,...,5} {%
    \ifnum\dsj<3\def\dscf{udsblue!28}\def\dscl{udsblue}\else\def\dscf{udsdark!20}\def\dscl{udsdark}\fi
    \filldraw[fill=\dscf,draw=\dscl,line width=0.3pt] ({#1+#3*\dsj},#2) rectangle ++(#3,-#4);}}
% \codematrix{x}{y}{columns}{rows}{cell}: latent codes as a grid of blue cells of random shade, top-left at
% (x,y); a fixed seed keeps the picture the same in every build
\newcommand\codematrix[5]{%
  \pgfmathsetseed{625}%
  \foreach \cmr in {1,...,#4} \foreach \cmc in {1,...,#3} {%
    \pgfmathtruncatemacro\cmshade{15+55*rnd}%
    \fill[udsblue!\cmshade] ({#1+(\cmc-1)*#5},{#2-(\cmr-1)*#5}) rectangle ++(#5,-#5);}
  \draw[udsblue,line width=0.5pt] (#1,#2) rectangle ({#1+#3*#5},{#2-#4*#5});}
% \minitree{x}{y}: a regression tree of depth 2, its root at (x,y), 7 mm wide and 6 mm tall
\newcommand\minitree[2]{%
  \draw[udsgreymid,line width=0.4pt] (#1,#2) -- (#1-2,#2-2.8) (#1,#2) -- (#1+2,#2-2.8)
    (#1-2,#2-2.8) -- (#1-3,#2-5.4) (#1-2,#2-2.8) -- (#1-1,#2-5.4)
    (#1+2,#2-2.8) -- (#1+1,#2-5.4) (#1+2,#2-2.8) -- (#1+3,#2-5.4);
  \foreach \mtx/\mty in {0/0,-2/-2.8,2/-2.8} \filldraw[fill=bluetint,draw=udsblue,line width=0.4pt] (#1+\mtx,#2+\mty) circle (0.75);
  \foreach \mtx in {-3,-1,1,3} \filldraw[fill=white,draw=udsgreymid,line width=0.4pt] (#1+\mtx-0.55,#2-4.85) rectangle ++(1.1,-1.1);}
